\documentclass[svgnames,x11names,dvipsnames]{beamer}
% --- Self-contained preamble (same conventions as the course decks).
% --- If you prefer, replace this block with \input{preamble}.
\usepackage[utf8]{inputenc}
\usepackage[english]{babel}
\usepackage{amsmath,amssymb}
\usepackage{xcolor}
\usepackage{graphicx}
\usepackage{hyperref}
\usepackage{booktabs}
\usepackage{array}
\IfFileExists{lmodern.sty}{\usepackage{lmodern}}{}
\usepackage{tikz}
\usetikzlibrary{positioning,arrows.meta,calc,fit,decorations.pathreplacing}
\usepackage[font=footnotesize]{caption}
\usepackage{qrcode}

\definecolor{color1bg}{HTML}{005e9e}
\setbeamercolor*{palette primary}{fg=white,bg=color1bg}
\setbeamercolor*{block title}{bg=color1bg,fg=white}
\setbeamercolor*{frametitle}{fg=color1bg}
\setbeamertemplate{caption}[numbered]
\definecolor{color2bg}{HTML}{003a63}
\setbeamercolor{block title alerted}{bg=color2bg,fg=white}
\setbeamercolor{block body alerted}{bg=color2bg!8}
\setbeamercolor{alerted text}{fg=color2bg}
\beamertemplatenavigationsymbolsempty

% --- This deck: one blue everywhere, the blue of the title page (Beamer's default, #3333B2).
% --- Delete this block to go back to color1bg / color2bg.
\definecolor{titleblue}{HTML}{3333B2}
\definecolor{accentred}{HTML}{B22222}
\colorlet{color1bg}{titleblue}
\colorlet{color2bg}{titleblue!75!black}
\setbeamercolor*{frametitle}{fg=titleblue}
\setbeamercolor*{block title}{bg=titleblue,fg=white}
\setbeamercolor{block title alerted}{bg=titleblue,fg=white}
\setbeamercolor{block body alerted}{bg=titleblue!8}
\setbeamercolor{alerted text}{fg=titleblue}

% --- Change marker: everything modified since the 15 October draft is in ... (bold).
% --- New or rewritten frames have their title in \chg; inside kept frames, only the changed text.
% --- Switched off in this version; comment the second line to see the changes.
\newcommand{\chg}[1]{{\bfseries #1}}
\renewcommand{\chg}[1]{#1}

\title[LLMs for Social Science Research]{\large \textbf{Large Language Models for Social Science Research} \\
    \textit{Session 2. Structured Data with an Agent}}
\author[H. Subtil]{\textbf{Hugo Subtil}}
\institute{PhD Workshop, University of Zurich}
\date{October 28, 2026}

\begin{document}
\frame{\titlepage}

% ============================================================================
% SESSION OVERVIEW (first, as requested)
% ============================================================================

\begin{frame}{Today's Session Overview (9:00--12:00)}

\setlength{\leftmargini}{2.3cm}
\small
\begin{enumerate}
    \setlength{\itemsep}{1pt}
    \item[\textbf{9:00--9:10}] Last Week, and What to Keep in Mind {\footnotesize (10 min)}

    \item[\textbf{9:10--9:30}] How I Actually Work {\footnotesize (20 min)}

    \item[\textbf{9:30--9:45}] Tokens and Models {\footnotesize (15 min)}

    \item[\textbf{9:45--10:25}] Block 1. Specify {\footnotesize (40 min)}, then a 15-min break
    \begin{itemize}
        \item Your analysis plan, with the agent (06)
    \end{itemize}

    \item[\textbf{10:40--11:15}] Block 2. Implement {\footnotesize (35 min)}
    \begin{itemize}
        \item The agent codes, you read the row counts (07)
    \end{itemize}

    \item[\textbf{11:15--11:45}] Block 3. Review and Decide {\footnotesize (30 min)}
    \begin{itemize}
        \item A fresh session reviews, you decide (08, then 05)
    \end{itemize}

    \item[\textbf{11:45--12:00}] Debrief and Next Session {\footnotesize (15 min)}
\end{enumerate}

\vspace{0.25cm}
{\footnotesize Slower than last time: every block is \emph{I show} (10 min), \emph{you do}, \emph{we check} (5 min).\\
Today is structured data. Brought a text corpus? Keep it for Sessions 3 and 4.}

\end{frame}

% ============================================================================
% PART 1: LAST WEEK
% ============================================================================

\begin{frame}[plain]
  \centering
  \vfill
  {\LARGE \textbf{\textcolor{titleblue}{Last Week}}}\\[0.4cm]
  {\large \textit{9:00 -- 9:10}}\\[0.3cm]
  {What we saw, and what to keep in mind}
  \vfill
\end{frame}

\begin{frame}{Last Week, in One Slide}

\small
\textbf{The argument.} An AI system now writes economics papers on its own (APE). Generating a paper is almost free; verifying one is not. The bottleneck moved from \emph{doing} to \emph{judging}, and the chain has to end with a human who understands.

\vspace{0.25cm}

\textbf{The position.} Delegate the work. Keep the skill to check it, and the responsibility for it. Three rules:
\begin{enumerate}
    \setlength{\itemsep}{0pt}
    \item You delegate execution, never understanding.
    \item Read before you accept.
    \item Stop when you no longer know what it is doing.
\end{enumerate}

\vspace{0.25cm}

\textbf{The tool.} A folder that remembers: \texttt{README.md} (what the project is), \texttt{CLAUDE.md} (how the agent behaves), \texttt{PROJECT\_HISTORY.md} (where we stand). Raw data read-only, prompts versioned, every session ends with a commit.

\vspace{0.25cm}

\textbf{The exercise.} Literature notes on your papers, then one note checked against the PDF, looking for the agent's mistake.

\end{frame}

\begin{frame}{What to Keep in Mind}

\textbf{This is still an open debate.} How much to delegate, and how much to check, is not settled. Reasonable people draw the line in different places.

\vspace{0.45cm}

\textbf{Since last week, I asked many colleagues one question:}

\vspace{0.3cm}

\centering
{\Large \alert{Where is your limit on review?}}\\[0.2cm]
{\small What do you read line by line, and what do you let go?}

\vspace{0.45cm}
\raggedright

\textbf{Today, one answer I find convincing}, from Germain Gauthier (Bocconi).
Not \emph{the} answer: a workable one, which you can adapt.

\vspace{0.35cm}

\end{frame}

% ============================================================================
% PART 2: HOW I ACTUALLY WORK (after Germain Gauthier)
% ============================================================================

\begin{frame}[plain]
  \centering
  \vfill
  {\LARGE \textbf{\textcolor{titleblue}{How I Actually Work}}}\\[0.4cm]
  {\large \textit{9:10 -- 9:30}}\\[0.3cm]
  {Inspired by Germain Gauthier (Bocconi)}
  \vfill
\end{frame}

\begin{frame}{How I Use AI}

\begin{itemize}
    \item AI is part of pretty much everything I do on my laptop.
\end{itemize}

\vspace{0.2cm}

\textbf{Some use cases}
\begin{itemize}
    \setlength{\itemsep}{3pt}
    \item A quick overview of papers I am unfamiliar with
    \item Writing code, organising project folders, updating results in a paper
    \item Proof-checking that the code, the paper and the slides are consistent
    \item Discussing ideas and analyses while I work on a project
    \item Improving the clarity of my writing and presentations
\end{itemize}

\vspace{0.4cm}

\textbf{Pen and paper} when I need to fully understand something (theory, statistics), or for mental clarity (a to-do list, planning ahead).

\end{frame}

\begin{frame}{How I Do Not Use AI}

\textbf{The one thing not to do:}

\vspace{0.15cm}
\centering
\textit{``Write me a full paper and a full slide set on this vague idea\\ \chg{while I go swim in the Limmat.}''}

\vspace{0.4cm}
\raggedright

\begin{itemize}
    \setlength{\itemsep}{6pt}
    \item \textbf{First, current AIs are terrible at this when left unattended.} It is a very iterative process, and you need to be in the loop.
    \item \textbf{Second, you will present the paper and the results.} Being in the loop is how you understand what you are doing.
\end{itemize}

\vspace{0.3cm}

\centering
\alert{The journey is almost as important as the destination in research.}

\vspace{0.3cm}
\raggedright
{\small AI can still speed things up, and by a lot. Last week's APE is what happens when nobody is in the loop.}

\end{frame}

\begin{frame}{The Setup: One Project Folder}

\small
\begin{columns}[T]
\column{0.52\textwidth}
\textbf{Germain's agent is connected to five tools}

\vspace{0.1cm}
{\footnotesize
\begin{tabular}{@{}ll@{}}
Dropbox & project files and data \\
GitHub & code history and repository \\
Overleaf & paper and slides \\
HPC cluster & heavy computation \\
Gmail & co-authors' comments \\
\end{tabular}}

\vspace{0.35cm}
\textbf{The project folder}

\vspace{0.1cm}
{\footnotesize
\begin{tabular}{@{}ll@{}}
\texttt{my-project/} & \\
\texttt{\ \ overleaf/} & local copy of the paper \\
\texttt{\ \ code/} & code repository \\
\texttt{\ \ data/} & raw and processed data \\
\texttt{\ \ notes/} & working notes \\
\end{tabular}}

\column{0.45\textwidth}
\textbf{How it runs}
\begin{itemize}
    \setlength{\itemsep}{3pt}
    \item The agent runs from the project folder.
    \item He talks to it in the browser, to see PDFs and figures as they are produced.
    \item The local copy of the paper and the Overleaf project are always in sync.
\end{itemize}
\end{columns}

\vspace{0.45cm}

\centering
\normalsize
The same idea as your folder from last week.\\
You do not need five tools today: one git folder is enough.

\end{frame}

\begin{frame}{Principle 1: You Are in Charge}

\begin{itemize}
    \setlength{\itemsep}{8pt}
    \item Agents overproduce, expand the scope, and overreact to details.
    \item Do not let them run unattended for long (overnight), unless the task is very well defined: for example, monitoring jobs on a cluster and relaunching the ones that fail.
    \item Stay hands-on: in the loop, correcting course, making the judgement calls.
    \item Being hands-on is also how you keep understanding what you are doing, and keep a map of it in your head, even when you no longer write the code.
\end{itemize}

\end{frame}

\begin{frame}{Principle 2: Accuracy, Not Speed or Cost-Cutting}

\begin{itemize}
    \setlength{\itemsep}{8pt}
    \item Run the models at maximum reasoning. For \emph{everything}.
    \item A mistake in a paper costs far more than extra tokens or waiting time.
    \item The difference in intelligence between reasoning levels is large.
    \item More reasoning does not remove the need for verification.
\end{itemize}

\vspace{0.6cm}

\centering
You asked how to save tokens. This principle says: not here.\\
At 9:30: what it costs, and where you can save without breaking it.

\end{frame}

\begin{frame}{A Typical Loop}



\small
\begin{enumerate}
    \setlength{\itemsep}{6pt}
    \item \textbf{Specify the task in writing.} Sample, equation, estimator, clustering, expected output: as in a detailed pre-analysis plan. Germain writes it in Overleaf; you, in \texttt{notes/analysis\_plan.md}.
    \item \textbf{A first agent implements it.} It flags ambiguities \emph{before} coding: there is high value in being on the same page.
    \item \textbf{Other agents review it in a fresh context.} They read the plan, the code and the outputs, and list the discrepancies.
    \item \textbf{You decide what to revise or accept.} A reviewer asked to find problems always finds something. Resolve, rerun, repeat until satisfied.
\end{enumerate}

\vspace{0.2cm}
{\footnotesize Not exotic: Anthropic's own guide to Claude Code recommends the same steps.}

\end{frame}

\begin{frame}{\chg{The Typical Failure: Silent Decisions}}

\chg{A typical failure mode is the AI making consequential decisions without asking. For example, dropping some observations.}

\vspace{0.3cm}

\textbf{On data, it looks like this}
\begin{itemize}
    \setlength{\itemsep}{3pt}
    \item rows lost in a merge (inner join instead of left join)
    \item a missing code (\texttt{-9}, \texttt{99}) read as a value
    \item a duplicate removed
    \item an identifier read as a number (\texttt{01001} becomes \texttt{1001})
\end{itemize}

\vspace{0.3cm}

\textbf{Two defences}
\begin{itemize}
    \setlength{\itemsep}{3pt}
    \item Every script prints the number of rows before and after each step. A change the plan does not explain stops the agent.
    \item I always ask the reviewing agents to check for this too.
\end{itemize}

\vspace{0.3cm}

\centering
Last week's question: how would you know a merge dropped a tenth of your observations?

\end{frame}

\begin{frame}{\chg{Two Phases: Prototyping, Then Sharpening}}

\textbf{Prototyping phase:}
\begin{itemize}
    \setlength{\itemsep}{3pt}
    \item It is an exploratory phase.
    \item I do not review the code or the prose.
    \item I treat these outputs as probable but provisional.
\end{itemize}

\vspace{0.4cm}

\textbf{Sharpening phase:}
\begin{itemize}
    \setlength{\itemsep}{3pt}
    \item If I think there is something there, I get involved in the writing and coding process.
    \item I audit the code, paying special attention to the number of rows in the dataset and results that ``look weird''. I have many independent agents check the code.
    \item I rewrite the paper and work with AI to iterate on the visuals.
    \item I take responsibility for the results and claims.
\end{itemize}

\end{frame}

\begin{frame}{Communication with Co-authors}

\small
Keep three distinctions explicit.

\vspace{0.15cm}
{\footnotesize
\begin{tabular}{@{}ll@{}}
\toprule
\textbf{Origin} & what the AI produced; what I produced \\
\textbf{Status} & what is exploratory; what is polished \\
\textbf{Responsibility} & what I have not checked; what I am ready to stand behind \\
\bottomrule
\end{tabular}}

\vspace{0.3cm}

Your co-authors will not look at exploratory and polished results in the same way, and will not make the same comments. AI use and understanding vary across co-authors: be transparent, and take responsibility for your work.

\vspace{0.3cm}

\end{frame}

% ============================================================================
% PART 3: TOKENS AND MODELS
% ============================================================================

\begin{frame}[plain]
  \centering
  \vfill
  {\LARGE \textbf{\textcolor{titleblue}{Tokens and Models}}}\\[0.4cm]
  {\large \textit{9:30 -- 9:45}}\\[0.3cm]
  {What Principle 2 costs, and where you can save}
  \vfill
\end{frame}

\begin{frame}{\chg{What Costs Tokens}}

\begin{itemize}
    \setlength{\itemsep}{8pt}
    \item \textbf{Every message re-sends the whole conversation} to the model: the files it read, its outputs, all previous turns. Turn 40 re-reads turns 1 to 39.
    \item \textbf{The cache discounts what was already sent}, while it is warm: about an hour on a subscription, five minutes with an API key. Back from lunch, the first message pays full price.
    \item \textbf{Effort multiplies tokens}: more files read, more checks, more thinking.
    \item \textbf{A bigger model costs more per turn.}
    \item \textbf{Subagents} spend tokens while you are not typing.
\end{itemize}

\vspace{0.4cm}

\centering
Subscription: tokens are quota. API key: tokens are money.\\
\texttt{/usage} shows where yours went.

\end{frame}

\begin{frame}{\chg{Model and Effort: Two Dials}}

\textbf{Model: how capable.} A bigger model understands more: unfamiliar libraries, ambiguous instructions, subtle bugs. It costs more per token.

\vspace{0.3cm}

\textbf{Effort: how thorough.} The same model works longer: it reads more files, thinks more, checks its own work before answering. It costs more tokens.

\vspace{0.5cm}

\textbf{When the answer is wrong, ask which one it is.}
\begin{itemize}
    \setlength{\itemsep}{4pt}
    \item It \emph{did not know}: it misread an ambiguous variable, or did not know how a package behaves. Change the \textbf{model}.
    \item It \emph{did not try}: it skipped a file, did not rerun the code, did not check. Raise the \textbf{effort}.
\end{itemize}

\vspace{0.4cm}

\centering
\texttt{/model} to see and change the model. \texttt{/effort} for the effort:\\
\texttt{low}, \texttt{medium}, \texttt{high}, \texttt{xhigh}, \texttt{max}.

\end{frame}

\begin{frame}{\chg{Which Model, Which Effort, for Which Task}}

\centering
{\footnotesize
\renewcommand{\arraystretch}{1.25}
\setlength{\tabcolsep}{5pt}
\begin{tabular}{@{}>{\raggedright\arraybackslash}p{4.1cm}ll>{\raggedright\arraybackslash}p{3.3cm}@{}}
\toprule
\textbf{Task} & \textbf{Model} & \textbf{Effort} & \textbf{Why} \\
\midrule
Rename, reformat a table, commit message & Sonnet & \texttt{low} & you check each result anyway \\
Prototyping: exploratory scripts, quick plots & Sonnet & \texttt{medium} & provisional by design \\
Writing: the paper, memos, the analysis plan & \chg{Fable} & \chg{\texttt{high}} & \chg{long, ambiguous tasks} \\
Code that enters the paper: cleaning, merging, estimation & Opus & \texttt{high} & errors hide in cases that are easy to overlook \\
Independent review & Opus & \texttt{xhigh} & the last check before yours \\
\bottomrule
\end{tabular}}

\vspace{0.5cm}

{\footnotesize Default in Claude Code: Opus 5.5 at \texttt{medium}. Fable is the most expensive per token: select it with \texttt{/model fable}, and keep it for what needs it.}

\end{frame}

\begin{frame}{Save on Context, Not on Reasoning}

\small
\begin{itemize}
    \setlength{\itemsep}{2pt}
    \item \textbf{One task, one session.} \texttt{/clear} between unrelated tasks costs nothing and loses nothing: the project's memory is in the folder (\texttt{PROJECT\_HISTORY.md}), not in the chat.
    \item \textbf{Watch the first minute.} \texttt{Esc} as soon as it drifts; \texttt{/rewind} to go back. After two failed corrections: \texttt{/clear} and one better prompt, not a third correction.
    \item \textbf{Never paste data.} The agent writes a script that prints summaries (types, counts, missing values); only the summaries enter the context. It also protects confidential data.
    \item \textbf{Point, do not paste.} \texttt{@code/03\_merge.py}, not the file's content. PDFs: only the pages you need.
    \item \textbf{A short \texttt{CLAUDE.md}}: under about 200 lines. Long procedures go in \texttt{prompts/}.
    \item \textbf{A long task:} \texttt{/compact} with a focus (``keep the merge decisions''). Back after more than an hour: a fresh session costs less.
\end{itemize}

\vspace{0.15cm}

\centering
\normalsize
Save on context. Never on reasoning for what enters the paper.

\end{frame}

% ============================================================================
% WHAT YOU LEAVE WITH
% ============================================================================

\begin{frame}[plain]
  \centering
  \vfill
  {\LARGE \textbf{\textcolor{titleblue}{\chg{Your Project, Three Blocks}}}}\\[0.4cm]
  {\large \textit{\chg{9:45 -- 11:45}}}\\[0.3cm]
  {\chg{I show, you do, we check}}
  \vfill
\end{frame}

\begin{frame}{What You Leave With at Noon}

In your folder:
\begin{itemize}
    \item \texttt{notes/analysis\_plan.md}, in your words, with its open questions settled
    \item numbered scripts in \texttt{code/} that print their row counts
    \item a cleaned dataset in \texttt{data/processed/} and a first table or figure in \texttt{outputs/}
    \item a review written by a fresh session, and your decision on each point it raised
    \item an entry in \texttt{PROJECT\_HISTORY.md} \chg{and a commit}
\end{itemize}

\vspace{0.3cm}

\textbf{And one thing you will have seen:} a decision the agent took on your data without asking. Finding it is Block 3. \alert{It is the only step you cannot skip.}

\end{frame}

% ============================================================================
% PART 4: THREE BLOCKS
% ============================================================================

\begin{frame}{Block 1. Specify (9:45--10:25)}

\small
\textbf{I show (10 min).} The example project: two raw files, election results and municipal indicators. Prompt 06: the agent inspects the files without changing them and interviews me. I answer; it writes the plan.

\vspace{0.3cm}

\textbf{You do (25 min)}
\begin{enumerate}
    \setlength{\itemsep}{3pt}
    \item Your data in \texttt{data/raw/}, its source in \texttt{data/raw/README.md}. No data? The example pack, in the course repository.
    \item A new session in your project folder. Paste prompt 06, from the course repository (\texttt{template/prompts/}).
    \item Answer its questions. What you do not know yet goes to ``Open questions''. It asks; it does not decide.
    \item Read the plan, then let it save \texttt{notes/analysis\_plan.md}.
\end{enumerate}

\vspace{0.3cm}

\textbf{We check (5 min).} Two plans on screen. What is in their ``Open questions''? Which merge has no rule for the rows that do not match?

\end{frame}

\begin{frame}{Block 2. Implement (10:40--11:15)}

\small
\textbf{I show (10 min).} \texttt{/clear}, then prompt 07. The agent lists what the plan leaves open; I answer. It writes the scripts and prints the row counts. We read them together: where did the rows go?

\vspace{0.3cm}

\textbf{You do (20 min)}
\begin{enumerate}
    \setlength{\itemsep}{3pt}
    \item \texttt{/clear}: the plan is in the folder, not in the chat. Paste prompt 07.
    \item Answer its list of ambiguities before it writes any code.
    \item Read the row-count trail it gives you, step by step.
    \item Rows that change without the plan saying so: paste the step in the shared document.
\end{enumerate}

\vspace{0.3cm}

\textbf{We check (5 min).} The shared document: who lost rows, at which step, and was it in the plan?

\end{frame}

\begin{frame}{Block 3. Review and Decide (11:15--11:45)}

\small
\textbf{I show (10 min).} I close the session and open a fresh one. Prompt 08: a reviewer that did not see the code being written. I go through its review and decide, point by point.

\vspace{0.3cm}

\textbf{You do (15 min)}
\begin{enumerate}
    \setlength{\itemsep}{3pt}
    \item Close the session. Open a new one. Paste prompt 08. \alert{Not optional.}
    \item For each point of the review: fix it, or accept it and write down why.
    \item Rerun. Every change in the row counts now follows from the plan.
    \item Prompt 05: the entry in \texttt{PROJECT\_HISTORY.md}, \chg{and} the commit.
\end{enumerate}

\vspace{0.3cm}

\textbf{We check.} That is the debrief.

\end{frame}

% ============================================================================
% WRAP-UP
% ============================================================================

\begin{frame}{Debrief}

\textbf{Three screens.} An analysis plan. A row-count trail. A review, with the discrepancy it found.

\vspace{0.3cm}

\textbf{The decisions the agent took alone today.} Read from the shared document. The list is the lesson.

\vspace{0.3cm}

\textbf{One question.} The table you would have sent your supervisor an hour ago: which \chg{status, exploratory or polished,} would it have carried?

\vspace{0.3cm}

\textbf{Before you leave.} \chg{Ask the agent to copy the section ``Row counts and silent decisions'' of \texttt{template/CLAUDE.md} into yours. Then add one line of your own that would have caught today's error.}

\end{frame}

\begin{frame}{Next Session: 4 November 2026}

\textbf{Session 3: Building Datasets from APIs, and Classic Text-as-Data}

\vspace{0.2cm}

\begin{itemize}
    \item Collecting structured and unstructured data through APIs, with the agent writing the collector
    \item Information extraction without an LLM: dictionaries, embeddings
    \item Bring one source you would like to collect from: an API, a website, a set of documents
    \item Keep working in your folder. End every session with prompt 05.
\end{itemize}

\vspace{0.3cm}

\begin{alertblock}{Carry-over question}
An API returns 10,000 documents and the agent reports that the collection is complete. How would you know?
\end{alertblock}

\end{frame}

% ============================================================================
% REFERENCES
% ============================================================================

\begin{frame}{References}

{\scriptsize

Gauthier, G.\ (2026). Using AI in Research Workflows. Slides, Bocconi University, September 2026. \href{https://pinchofdata.github.io/germaingauthier/pdfs/teaching/ai_research_workflow.pdf}{\textcolor{titleblue}{pinchofdata.github.io/germaingauthier}}. The section ``How I Actually Work'' follows his slides closely.

\vspace{0.1cm}

Anthropic (2026). Best Practices for Claude Code. \href{https://code.claude.com/docs/en/best-practices}{\textcolor{titleblue}{code.claude.com/docs/en/best-practices}}.

\vspace{0.1cm}

Anthropic (2026). Model Configuration: Aliases and Effort Levels. \href{https://code.claude.com/docs/en/model-config}{\textcolor{titleblue}{code.claude.com/docs/en/model-config}}.

\vspace{0.1cm}

Anthropic Support (2026). Models and Usage Limits in Claude Code. \href{https://support.claude.com/en/articles/14552983-models-usage-limits-in-claude-code}{\textcolor{titleblue}{support.claude.com}}.

\vspace{0.1cm}

Hallie, L.\ (2026). Choosing a Claude Model and Effort Level in Claude Code. \textit{Claude Blog}, 7 July 2026.

}

\end{frame}

\end{document}